\subsection{Solution to Problem 25: Complex Numbers with Three Conditions}

\begin{solution}
Since $|z_1| = |z_2| = |z_3| = r$, write
\[
z_k = r(\cos\theta_k + i\sin\theta_k) \qquad (k = 1,2,3).
\]
From $z_1 z_2 z_3 = 1$,
\[
r^3\bigl[\cos(\theta_1 + \theta_2 + \theta_3) + i\sin(\theta_1 + \theta_2 + \theta_3)\bigr] = 1.
\]
The right-hand side is the real number $1$, so the imaginary part is zero and the modulus is $1$. Hence $r^3 = 1$ (so $r = 1$) and
\[
\theta_1 + \theta_2 + \theta_3 = 2\pi k \quad \text{for some integer } k.
\]
Thus each $z_k$ lies on the unit circle: $z_k = \cos\theta_k + i\sin\theta_k$.

From $z_1 + z_2 + z_3 = 1$,
\[
(\cos\theta_1 + \cos\theta_2 + \cos\theta_3) + i(\sin\theta_1 + \sin\theta_2 + \sin\theta_3) = 1.
\]
The imaginary parts must sum to zero. A natural symmetric choice is to take one number real and the other two a conjugate pair: let
\[
z_1 = 1, \qquad z_2 = \cos\theta + i\sin\theta, \qquad z_3 = \cos\theta - i\sin\theta.
\]
Then
\[
z_1 + z_2 + z_3 = 1 + 2\cos\theta = 1
\quad\Longrightarrow\quad
\cos\theta = 0,
\]
so $\theta = \dfrac{\pi}{2}$ or $\dfrac{3\pi}{2}$.

For $\theta = \dfrac{\pi}{2}$: $z_1 = 1$, $z_2 = i$, $z_3 = -i$.

Check: $|1| = |i| = |-i| = 1$ \checkmark, $1 + i + (-i) = 1$ \checkmark, $1 \cdot i \cdot (-i) = 1$ \checkmark.

For $\theta = \dfrac{3\pi}{2}$: $z_2 = -i$, $z_3 = i$, giving the same set permuted.

By symmetry, any of $z_1$, $z_2$, $z_3$ may be the real number $1$, with the other two equal to $i$ and $-i$.

Therefore the solutions are all permutations of $\{1, i, -i\}$.
\end{solution}

\begin{takeaways}
\begin{itemize}
    \item Equal moduli: Write each number in modulus--argument form with a common radius $r$
    \item Product condition: $z_1z_2z_3 = 1$ forces $r = 1$ and a sum of arguments equal to $2\pi k$
    \item Symmetry: A conjugate pair together with a real number simplifies the sum condition
\end{itemize}
\end{takeaways}
